\section{Arithmetic quartet parity and the exact positive frontier operator}
\label{sec:arithmetic-frontier-parity}

Let \(h\) be a nonempty packet of actual zeros with their full
orders, closed under conjugation and under
\(\rho\mapsto1-\overline\rho\).  Their composition also
preserves the packet.  Retain \(d=\deg h\), \(g=2\xi\),
\(v_h=g/h\), \(\upsilon=j_hv_h\), and the actual measure
\(d\nu_h=|v_h(1/2+it)|^2dt/(2\pi)\).  Throughout this section
\(k\ge1\), \(M\ge k(d-1)\), and
\[
 E_k=(\mathbb C[s]/h)^{\otimes k},\qquad
 A_k=\sum_{j=1}^k A_{h,j},\qquad
 \mathcal C_M=G_M^{-1},\qquad
 W_M=A_k^*G_M+G_MA_k-kG_M.
 \tag{AP.1}
\]
The metric \(G_M\) is the actual total-degree minimum of
(CT.12)--(CT.13).  All coordinates below remain its original
monomial tensor coordinates.

\subsection{The actual two frontier matrices}

Let \(q_j\) be the original monic orthogonal polynomials for
\(\nu_h\), and \(\kappa_j=\int|q_j|^2d\nu_h\).
Thus \(\kappa_0\) retains the mass of the original measure.
For a multi-index \(\alpha\), put
\[
 q_\alpha=\prod_{j=1}^kq_{\alpha_j}(s_j),\qquad
 \kappa_\alpha=\prod_{j=1}^k\kappa_{\alpha_j},\qquad
 z_\alpha=U_{\upsilon^{\otimes k}}[q_\alpha]\in E_k.
 \tag{AP.2}
\]
The source's notation \(p_j,\omega_j\) denotes these same
\(q_j,\kappa_j\): their monic leading terms and identical
lower-degree orthogonality determine the polynomial uniquely,
and then determine its norm.  No polynomial or measure is divided
by its norm.  The triangular polynomial change of basis and product
integration in (CT.12) give exactly
\[
 \mathcal C_M=\sum_{|\alpha|\le M}
                       \frac{z_\alpha z_\alpha^*}{\kappa_\alpha}.
 \tag{AP.3}
\]
Indeed this basis has Gram diagonal \(\kappa_\alpha\),
and its full interpolation matrix has columns \(z_\alpha\);
substitution in \(JM^{-1}J^*\) proves (AP.3), retaining
\(\upsilon^{\otimes k}\) in every column.

Write \(a_j=\kappa_j/\kappa_{j-1}\) for \(j\ge1\),
\(a_0=0\), and retain the ordered frontier \(|\beta|=M+1\).
Its size is \(r=\binom{M+k}{k-1}\).  Define maps with the
same declared source \(\mathbb C^r\):
\[
 \begin{gathered}
 \Omega=\operatorname{diag}(\kappa_\beta)_{|\beta|=M+1},\qquad
 F=[z_\beta]_{|\beta|=M+1},\\
 w_\beta=\sum_{j:\beta_j>0}a_{\beta_j}z_{\beta-e_j},\qquad
 E_+=[w_\beta]_{|\beta|=M+1}.
 \end{gathered}
 \tag{AP.4}
\]
The original recurrence
\(sq_j=q_{j+1}+b_jq_j-a_jq_{j-1}\), with
\(b_j+\overline b_j=1\), applied to (AP.3) gives
\[
 \mathcal Z_M:=A_k\mathcal C_M+\mathcal C_MA_k^*-k\mathcal C_M
       =F\Omega^{-1}E_+^*+E_+\Omega^{-1}F^*,
 \qquad W_M=G_M\mathcal Z_MG_M.
 \tag{AP.5}
\]
To check the cancellation completely, an adjacent pair
\(\alpha,\alpha+e_j\) both of degree at most \(M\)
has coefficients \(\kappa_\alpha^{-1}\) and
\(-a_{\alpha_j+1}\kappa_{\alpha+e_j}^{-1}\), whose
sum is zero.  Its transposed conjugate cancels in the same way.
The diagonal contributes
\(\sum_j(b_{\alpha_j}+\overline b_{\alpha_j})=k\),
which cancels the last term.  The remaining frontier pair has
\(\beta=\alpha+e_j\) and coefficient
\(\kappa_\alpha^{-1}=a_{\beta_j}/\kappa_\beta\);
summing these pairs is exactly (AP.4)--(AP.5).
The same sum gives the full inverse-metric update
\[
 \mathcal C_{M+1}=\mathcal C_M+F\Omega^{-1}F^*.
 \tag{AP.6}
\]

\subsection{The original linear reflection separates the two frontiers}

Conjugation closure makes \(h\) real, and reflection closure
then gives \(h(1-s)=(-1)^dh(s)\).  The original functional
equation gives
\[
 v_h(1-s)=(-1)^dv_h(s),\qquad
 q_j(1-s)=(-1)^jq_j(s).
 \tag{AP.7}
\]
For the second identity, the actual density is even in \(t\).
The polynomial \((-1)^jq_j(1-s)\) is monic, and substitution
\(t\mapsto-t\) verifies its orthogonality to every lower
degree.  The difference from \(q_j\) has degree less than
\(j\) and is orthogonal to itself, so positivity proves it is
zero.  This proves (AP.7) with the original leading coefficient.

Let \(\mathfrak c:E_k\to E_k\) be the complex-linear map
substituting \(s_j\mapsto1-s_j\) in every factor.  It is
well-defined because the ideal \((h(s_1),\ldots,h(s_k))\)
is preserved, and \(\mathfrak c^2=I\).  Let \(C\) denote
its matrix in the retained original basis.  The original function
map
\[
 (\mathfrak t_kf)(x_1,\ldots,x_k)
       =\left(\prod_jx_j^{-1}\right)f(x_1^{-1},\ldots,x_k^{-1})
 \tag{AP.8}
\]
is a linear unitary involution for \(d^kx\), by the change of
variables \(y_j=x_j^{-1}\).  Its Mellin transform replaces
\(s_j\) by \(1-s_j\).  Consequently, on an original numerator,
\[
 \mathfrak t_k\mathcal T_h^{(k)}P
     =(-1)^{kd}\mathcal T_h^{(k)}P(1-s_1,\ldots,1-s_k).
 \tag{AP.9}
\]
Equations (AP.7) and (AP.2) show that its full jet is
\(\mathfrak c\mathcal J P\), with exactly the displayed
unit sign.  Total degree and the original relation ideal are
preserved.  Applying \(\mathfrak t_k\) to the unique minimum
for a prescribed jet therefore gives the unique minimum for its
\(\mathfrak c\)-image.  Uniqueness follows from the orthogonal
decomposition (CT.5).  Hence
\[
 \mathfrak t_kR_M=R_MC,\qquad C^*G_MC=G_M,
 \qquad Cz_\alpha=(-1)^{kd+|\alpha|}z_\alpha.
 \tag{AP.10}
\]
In particular, with \(\eta=(-1)^{kd+M+1}\),
\[
 CF=\eta F,\qquad CE_+=-\eta E_+.
 \tag{AP.11}
\]
Every term defining \(w_\beta\) has degree \(M\), which
proves the second identity without suppressing any coefficient.

Declare the following maps between the original positive spaces
and their coordinate Hilbert spaces:
\[
 P=G_M^{1/2}CG_M^{-1/2},\qquad
 U=G_M^{1/2}F\Omega^{-1/2},\qquad
 V=G_M^{1/2}E_+\Omega^{-1/2}.
 \tag{AP.12}
\]
Here \(G_M^{1/2}\) and \(\Omega^{1/2}\) are the explicit
isometries from the declared original metric spaces to Euclidean
coordinate spaces, with their displayed inverses.  Every expression
in (AP.12) retains both original matrices.  From (AP.10), \(P\)
is unitary and \(P^2=I\); hence \(P^*=P\).  Equations
(AP.11) imply \(PU=\eta U\), \(PV=-\eta V\), and thus
\[
 U^*V=(PU)^*(PV)=-U^*V=0.
 \tag{AP.13}
\]

\subsection{The complete positive operator and the sharper arithmetic bound}

Put \(\mathsf A=U^*U\), \(\mathsf B=V^*V\), retaining their
source type \(\mathbb C^r\).  Equivalently,
\[
 \mathsf A=\Omega^{-1/2}F^*G_MF\Omega^{-1/2},\qquad
 \mathsf B=\Omega^{-1/2}E_+^*G_ME_+\Omega^{-1/2}.
 \tag{AP.14}
\]
These positive-semidefinite matrices are built from the two original
frontiers; \(\mathsf A\) is distinct from the arithmetic generator
\(A_k\).

\begin{theorem}[Exact positive frontier operator]
The actual relative weight satisfies
\[
 \begin{gathered}
 S_M=G_M^{-1/2}W_MG_M^{-1/2}=UV^*+VU^*,\\
 \boxed{\epsilon_{k,M}=\|UV^*\|,\qquad
 \epsilon_{k,M}^2=
       \lambda_{\max}(\mathsf A^{1/2}\mathsf B\mathsf A^{1/2}),}\qquad
 p_M=\operatorname{rank}(\mathsf A^{1/2}\mathsf B^{1/2}),\\
 \operatorname{Tr}(S_M^2)=2\operatorname{Tr}(\mathsf A\mathsf B).
 \end{gathered}
 \tag{AP.15}
\]
Its positive and negative inertias both equal \(p_M\).
Define the original incidence bound and volume ratio by
\[
 \Gamma_{k,M}=\max_{|\alpha|=M}
     \sum_{j=1}^k\bigl(a_{\alpha_j+1}+(k-1)a_{\alpha_j}\bigr),
 \quad \lambda_{k,M}=\|U\|^2,
 \quad \pi_M=\frac{\det G_{M+1}}{\det G_M}.
 \tag{AP.16}
\]
Then
\[
 \boxed{\epsilon_{k,M}^2
      \le\Gamma_{k,M}\lambda_{k,M}
      \le\Gamma_{k,M}(\pi_M^{-1}-1).}
 \tag{AP.17}
\]
All formulas include the zero-rank cases and every full arithmetic
jet.  No growth estimate for \(\Gamma_{k,M}\) or \(\pi_M\)
is assumed.
\end{theorem}
\begin{proof}
Equation (AP.5) proves the first equality in (AP.15).  Decompose
the retained coordinate Hilbert space into the two orthogonal
eigenspaces of the explicit Hermitian involution \(P\).
Relative to the order \((\eta,-\eta)\), the operator has blocks
\[
 S_M=\begin{pmatrix}0&T\\T^*&0\end{pmatrix},
 \qquad T=UV^*|_{\ker(P+\eta I)}.
 \tag{AP.18}
\]
Indeed \(U^*\) vanishes on the opposite eigenspace and \(V^*\)
vanishes on the first, by (AP.11).  Squaring gives diagonal
blocks \(TT^*,T^*T\).  For each nonzero singular value
\(\sigma\) of \(T\), choose unit singular vectors
\(x,y\) with \(Ty=\sigma x\), \(T^*x=\sigma y\);
then \((x,y)\) and \((x,-y)\) have eigenvalues
\(\sigma,-\sigma\).  The remaining eigenspaces are
\(\ker T^*\oplus\ker T\), with eigenvalue zero.  These
vectors form the full orthogonal decomposition by the spectral
theorem for \(T^*T\), including all its zero vectors.
It follows that \(\epsilon=\|T\|=\|UV^*\|\) and both
inertias equal \(\operatorname{rank}(UV^*)\).

Now \((UV^*)(UV^*)^*=U\mathsf B U^*\).  Its nonzero
eigenvalues and rank are those of
\(\mathsf B^{1/2}\mathsf A\mathsf B^{1/2}\): the maps
\(U\mathsf B^{1/2}\) and its adjoint transfer each nonzero
eigenvector, with inverse divided by that eigenvalue.  The latter
matrix is \(T_0^*T_0\), with
\(T_0=\mathsf A^{1/2}\mathsf B^{1/2}\), while
\(T_0T_0^*=\mathsf A^{1/2}\mathsf B\mathsf A^{1/2}\).
Their full positive spectra agree by the same maps, and their
ranks equal \(\operatorname{rank}T_0\).  This proves the
norm and rank formulas, including rank zero.  Finally (AP.13)
kills \((UV^*)^2\) and \((VU^*)^2\).  Taking the trace of
the remaining two products gives the last line of (AP.15).

For the incidence estimate, let \(Z=[z_\alpha]_{|\alpha|=M}\)
and \(\Omega_M=\operatorname{diag}(\kappa_\alpha)\).
Define \(L_{\alpha\beta}=a_{\beta_j}\) when
\(\alpha=\beta-e_j\), and zero otherwise; summation applies
if more than one specified incidence contributes.  Then \(E_+=ZL\).
The matrix \(N=\Omega_M L\Omega^{-1}L^*\) has nonnegative
entries.  Its row sum at \(\alpha\) is exactly
\[
 \sum_{j=1}^k\left(a_{\alpha_j+1}
                         +\sum_{l\ne j}a_{\alpha_l}\right)
   =\sum_{j=1}^k\bigl(a_{\alpha_j+1}+(k-1)a_{\alpha_j}\bigr).
 \tag{AP.19}
\]
To verify each factor, its \(j\)-th contributing frontier is
\(\beta=\alpha+e_j\), and
\(\kappa_\alpha a_{\beta_j}/\kappa_\beta=1\).
The predecessors of this \(\beta\) supply
\(\sum_l a_{\beta_l}\), which is precisely the displayed
summand.  Every eigenvalue of \(N\) has modulus at most the
largest row sum: in an eigenvector choose a coordinate of largest
absolute value and apply that row equation.  The matrix \(N\)
is similar, by \(\Omega_M^{1/2}\), to the Hermitian positive
matrix \(\Omega_M^{1/2}L\Omega^{-1}L^*\Omega_M^{1/2}\).
Its eigenvalues therefore lie in \([0,\Gamma_{k,M}]\).
Undoing the congruence and using (AP.3) proves
\[
 E_+\Omega^{-1}E_+^*
   \preceq\Gamma_{k,M}Z\Omega_M^{-1}Z^*
   \preceq\Gamma_{k,M}\mathcal C_M,
 \qquad VV^*\preceq\Gamma_{k,M}I.
 \tag{AP.20}
\]
Consequently \(\epsilon^2=\|UV^*\|^2
\le\|U\|^2\|V\|^2\le\Gamma_{k,M}\lambda_{k,M}\).
Finally (AP.6) gives
\[
 \pi_M^{-1}=\det(I+UU^*)
    =\prod_j(1+\lambda_j(UU^*))
    \ge1+\lambda_{k,M}.
 \tag{AP.21}
\]
All eigenvalues here are nonnegative, proving the second bound in
(AP.17).  If \(U=0\) or \(V=0\), (AP.15) gives zero
weight directly; all determinant denominators remain positive.
\end{proof}

The arithmetic departure has an exact expression in these same
two frontier Grams.  List the original \(\rho_i\) with every
algebraic multiplicity and put
\(\delta_i=\operatorname{Re}\rho_i-1/2\).  Comparing
(AP.15) with the fully proved tensor trace (CT.21) gives
\[
 \mathfrak d_M^{(k)}
    =\operatorname{Tr}(\mathsf A\mathsf B)
           -2\sum_{i_1,\ldots,i_k}
                         (\delta_{i_1}+\cdots+\delta_{i_k})^2.
 \tag{AP.22}
\]
This retains the complete nonsemisimple generator in the left-hand
definition of (CT.20).  For \(k=1\), a four-distinct-point
tesserine orbit of common full order \(\mu\) has the further
exact chain, with \(M=d+m\),
\[
 \mu^2|c_{-+}(\rho)|^2
 \le\epsilon_{1,M}^2
 =\lambda_{\max}(\mathsf A^{1/2}\mathsf B\mathsf A^{1/2})
 \le\Gamma_{1,M}(\pi_M^{-1}-1).
 \tag{AP.23}
\]
The first inequality, original character, multiplicity and critical
supported-zero case are proved in (R9); (AP.15)--(AP.17) prove
the rest.  Thus the original character observation is related to
an explicit positive frontier operator through the original
metric, full interpolation unit and declared isometries.

\subsection{The same positive frontier on the actual symmetric summand}

The cochain projector used here is exactly
\(\mathsf S_k=(k!)^{-1}\sum_{\pi\in S_k}
\operatorname{sgn}(\pi)T_\pi\), with the original Koszul
permutations \(T_\pi\).  Its top-degree action is the ordinary
permutation projection, because \(T_\pi\) there equals
\(\operatorname{sgn}(\pi)P_\pi\).  Equations
(SS.1)--(SS.9) prove its continuous cochain maps, completed
projective-tensor cohomology, and finite packet projector, including
the complementary summand.  We now use its finite invariant image
with the original total-degree representatives \(R_M\).

Write \(e_1,\ldots,e_d\) for the original power basis of
\(E\), with indices shifted by one, and retain each occupancy
\(n\in\mathbb N^d\), \(|n|=k\).  Let \(b_n\) be
the literal sum of the distinct basis tensors of occupancy \(n\).
Define the maps and dimensions by
\[
 \begin{gathered}
 N_s=\binom{k+d-1}{d-1},\qquad
 I:\mathbb C^{N_s}\longrightarrow E_k,\quad Ie_n=b_n,\\
 D_o=I^*I=\operatorname{diag}_{|n|=k}
                  \frac{k!}{\prod_i n_i!},\qquad
 L=D_o^{-1}I^*,\qquad P_s=IL=\frac1{k!}\sum_\pi P_\pi,
 \qquad LI=I_{N_s}.
 \end{gathered}                                                \tag{AP.24}
\]
Indeed the distinct original tensor basis vectors are orthonormal
for this coordinate calculation; hence their disjoint orbit sums
have precisely the displayed diagonal Gram \(D_o\).
The map \(L\) averages coefficients on each orbit, proving all
the compositions in (AP.24).  In particular \(P_s^*=P_s\)
and \(P_s^2=P_s\).  The factorials in \(D_o\) are
retained separately from the arithmetic Gram.

Permutation invariance of the original product measure, relation
ideal, total degree and full interpolation unit gives
\(R_MP_\pi=P_\pi R_M\), by applying the unitary inverse
permutation to each constrained minimum.  Thus \(G_M\) and
\(G_M^{-1}\) commute with \(P_s\).  Put
\[
 \begin{gathered}
 R_M^s=R_MI,\qquad G_M^s=I^*G_MI,\qquad W_M^s=I^*W_MI,\\
 A_k^s=LA_kI,\qquad A_kI=IA_k^s,\qquad
 W_M^s=(A_k^s)^*G_M^s+G_M^sA_k^s-kG_M^s,\\
 I^*G_M=G_M^sL,\qquad G_MI=L^*G_M^s,\qquad
 \mathcal K_M^s:=(G_M^s)^{-1}=LG_M^{-1}L^*.
 \end{gathered}                                                \tag{AP.25}
\]
These are the same representatives and matrices as (SS.10)--(SS.11).
For completeness, \(I^*G_M=I^*P_sG_M=I^*G_MP_s
=I^*G_MIL\), proving the first rectangular identity and then
its adjoint.  The generator commutes with every factor permutation,
so it preserves the image of \(I\), proving its intertwining
identity and the displayed weight formula.  Finally
\[
 (LG_M^{-1}L^*)(I^*G_MI)
 =LG_M^{-1}P_sG_MI=LI=I_{N_s},
\]
which proves the inverse formula with both factors \(D_o^{-1}\)
present.  Positivity follows since \(I\) is injective and
\(G_M\succ0\).  The notation \(\mathcal K_M^s\) refers
only to this inverse kernel; it will not denote a reflection.

Keep the complete, possibly redundant frontier of (AP.4), with its
unchanged domain \(\mathbb C^r\) and norm matrix \(\Omega\).
Its symmetric target maps are
\[
 \begin{gathered}
 F_s=LF:\mathbb C^r\longrightarrow\mathbb C^{N_s},\qquad
 E_s=LE_+:\mathbb C^r\longrightarrow\mathbb C^{N_s},\\
 \mathcal Z_M^s=F_s\Omega^{-1}E_s^*
                          +E_s\Omega^{-1}F_s^*,\qquad
 W_M^s=G_M^s\mathcal Z_M^sG_M^s,\\
 \mathcal K_{M+1}^s=\mathcal K_M^s
                                      +F_s\Omega^{-1}F_s^*.
 \end{gathered}                                                \tag{AP.26}
\]
The rectangular identities (AP.25), applied to
\(I^*G_M\mathcal Z_MG_MI\), prove the weight formula.
Applying \(L\) and \(L^*\) to (AP.6) proves the inverse
update.  It uses the same orbit maps at both source degrees.
Thus the update concerns the actual next symmetric minimum, not a
newly selected positive form.

The original linear reflection \(C\) commutes with every factor
permutation, since it substitutes the same \(1-s_j\) in each
factor.  Its invariant-coordinate matrix is
\[
 \begin{gathered}
 C_s=LCI,\qquad CI=IC_s,\qquad C_sL=LC,\qquad C_s^2=I_{N_s},\\
 C_s^*G_M^sC_s=G_M^s,\qquad
 C_sF_s=\eta F_s,\qquad C_sE_s=-\eta E_s,
 \quad\eta=(-1)^{kd+M+1}.
 \end{gathered}                                                \tag{AP.27}
\]
For example \(CI=P_sCI=ILCI\), and
\(C_sL=LCP_s=LP_sC=LC\); these prove the first
intertwiners and then the involution.  The metric equation follows
by substituting \(IC_s=CI\) into \(I^*G_MI\) and
using (AP.10).  Applying \(L\) to (AP.11) proves both
parities with the full arithmetic unit and its original sign.

Define the explicit isometries and the two retained frontier Grams by
\[
 \begin{gathered}
 P_M^s=(G_M^s)^{1/2}C_s(G_M^s)^{-1/2},\\
 U_s=(G_M^s)^{1/2}F_s\Omega^{-1/2},\qquad
 V_s=(G_M^s)^{1/2}E_s\Omega^{-1/2},\\
 \mathsf A_s=U_s^*U_s,\qquad \mathsf B_s=V_s^*V_s.
 \end{gathered}                                                \tag{AP.28}
\]
The superscript on \(P_M^s\) distinguishes this transported
reflection from the fixed permutation projection \(P_s\) in
(AP.24).  Equation (AP.27) proves that \(P_M^s\) is a
Hermitian involution and that \(U_s,V_s\) have opposite
eigenvalues \(\eta,-\eta\).  Their ranges are consequently
orthogonal, so \(U_s^*V_s=0\).

\begin{theorem}[Exact symmetric arithmetic frontier]
For the actual signed symmetric cohomology summand, the original
Hermitian pencil has
\[
 \begin{gathered}
 S_M^s=(G_M^s)^{-1/2}W_M^s(G_M^s)^{-1/2}
                           =U_sV_s^*+V_sU_s^*,\\
 \boxed{\epsilon_M^s=\|U_sV_s^*\|,\qquad
 (\epsilon_M^s)^2=
    \lambda_{\max}(\mathsf A_s^{1/2}\mathsf B_s
                                      \mathsf A_s^{1/2}),}\\
 p_M^s=\operatorname{rank}(\mathsf A_s^{1/2}\mathsf B_s^{1/2}),
 \qquad n_+(W_M^s)=n_-(W_M^s)=p_M^s,\\
 \operatorname{Tr}((S_M^s)^2)
                           =2\operatorname{Tr}(\mathsf A_s\mathsf B_s).
 \end{gathered}                                                \tag{AP.29}
\]
With the original \(\Gamma_{k,M}\) of (AP.16), put
\[
 \lambda_M^s=\|U_s\|^2,\qquad
 \pi_M^s=\frac{\det G_{M+1}^s}{\det G_M^s}.
\]
Then the same source sequence satisfies
\[
 \boxed{(\epsilon_M^s)^2\le\Gamma_{k,M}\lambda_M^s
          \le\Gamma_{k,M}((\pi_M^s)^{-1}-1),\qquad
          \lambda_M^s\le\lambda_{k,M}.}              \tag{AP.30}
\]
All formulas retain every occupancy and nilpotent direction in
(SS.12)--(SS.14), as well as the ungrouped frontier multiplicities.
\end{theorem}
\begin{proof}
Equation (AP.26) gives the first line of (AP.29).  Decomposing
\(\mathbb C^{N_s}\) into the two orthogonal eigenspaces of
\(P_M^s\) gives exactly the block operator (AP.18), now with
\(T_s=U_sV_s^*\) between those two spaces.  The singular-vector
maps and their inverses proved after (AP.18) apply to these declared
matrices without any rank assumption.  They give opposite eigenvalues
for each positive singular value, rank
\(\operatorname{rank}(\mathsf A_s^{1/2}\mathsf B_s^{1/2})\),
and all remaining eigenvalues zero.  The same two explicit composites
\(U_s\mathsf B_s^{1/2}\) and its adjoint identify the
positive eigenvalues with those of
\(\mathsf A_s^{1/2}\mathsf B_s\mathsf A_s^{1/2}\).
Finally \(U_s^*V_s=0\) kills the two square cross factors;
the traces of the other two products both equal
\(\operatorname{Tr}(\mathsf A_s\mathsf B_s)\).
This proves every assertion in (AP.29).

Applying \(L\) and \(L^*\) to the original inequality
(AP.20), and using (AP.25), proves
\[
 E_s\Omega^{-1}E_s^*
      \preceq\Gamma_{k,M}\mathcal K_M^s,\qquad
 V_sV_s^*\preceq\Gamma_{k,M}I_{N_s}.                  \tag{AP.31}
\]
Thus \(\|U_sV_s^*\|^2\le\|U_s\|^2\|V_s\|^2
\le\Gamma_{k,M}\lambda_M^s\).  Taking determinants of
the exact inverse update in (AP.26) gives
\[
 (\pi_M^s)^{-1}=\det(I_{N_s}+U_sU_s^*)
              =\det(I_r+\mathsf A_s)
              \ge1+\lambda_M^s,\qquad 0<\pi_M^s\le1.
                                                               \tag{AP.32}
\]
The determinant equality follows either by the two rectangular
composites and their nonzero spectra, or by multiplying the two
block factors for \(\det(I+XY)=\det(I+YX)\).
The inequality multiplies the nonnegative factors
\(1+\lambda_j(\mathsf A_s)\), exactly as in (AP.21).
This proves the first two bounds in (AP.30).

For the final bound, the literal orbit maps give
\[
 \begin{split}
 L^*G_M^sL&=P_sG_MP_s\preceq G_M,\\
 \mathsf A_s
 &=\Omega^{-1/2}F^*P_sG_MP_sF\Omega^{-1/2}
       \preceq\Omega^{-1/2}F^*G_MF\Omega^{-1/2}
        =\mathsf A.
 \end{split}                                                  \tag{AP.33}
\]
Indeed \(G_M\) commutes with the orthogonal projection
\(P_s\), so their omitted cross terms vanish and
\(G_M-P_sG_MP_s=(I-P_s)G_M(I-P_s)\succeq0\).
Taking the largest Rayleigh quotient of the two positive matrices
in (AP.33) proves \(\lambda_M^s\le\lambda_{k,M}\).
No comparison of arbitrary compressed inverse matrices is used.
\end{proof}

There is also an exact comparison of the two relative operators.
The map
\[
 Q_M=G_M^{1/2}I(G_M^s)^{-1/2}:
            \mathbb C^{N_s}\longrightarrow\mathbb C^{d^k}
\]
has \(Q_M^*Q_M=I_{N_s}\), and the inverse formula in
(AP.25) gives \(Q_MQ_M^*=P_s\).  Direct multiplication
therefore gives
\[
 S_M^s=Q_M^*S_MQ_M,\qquad
 U_s=Q_M^*U,\qquad V_s=Q_M^*V,\qquad
 \epsilon_M^s\le\epsilon_{k,M}.                       \tag{AP.34}
\]
Here \(G_M^{1/2}\) commutes with \(P_s\) by its
spectral decomposition.  The full \(S_M\) also commutes
with \(P_s\), since both \(A_k\) and \(G_M\) do.
Consequently (AP.34) is its invariant orthogonal restriction;
the complementary operator remains on \(\ker P_s\).
The norm inequality follows directly by applying \(S_M\)
to the unit vector \(Q_Mu\) and then the contraction
\(Q_M^*\).

If desired, the source frontier can be grouped with every factorial
retained.  For an orbit \(\mathcal O\) of degree-\(M+1\)
multi-indices, let
\(m_j(\beta)=|\{l:\beta_l=j\}|\) and
\[
 n_{\mathcal O}=\frac{k!}{\prod_{j\ge0}m_j(\beta)!},\qquad
 \widehat f_{\mathcal O}=\sum_{\beta\in\mathcal O}Lz_\beta,
 \quad\widehat e_{\mathcal O}=\sum_{\beta\in\mathcal O}Lw_\beta,
 \quad\widehat\omega_{\mathcal O}=n_{\mathcal O}\kappa_\beta.
                                                               \tag{AP.35}
\]
The product norm \(\kappa_\beta\) is constant on an orbit.
Equivariance of the full unit, polynomials and incidence sums implies
that \(Lz_\beta\) and \(Lw_\beta\) are individually
constant on the same orbit.  Hence, for example,
\[
 \sum_{\beta\in\mathcal O}
       \frac{(Lz_\beta)(Lw_\beta)^*}{\kappa_\beta}
       =\frac{\widehat f_{\mathcal O}
                    \widehat e_{\mathcal O}^*}
                    {\widehat\omega_{\mathcal O}}.     \tag{AP.36}
\]
The same calculation with both columns of the first or second type
proves the corresponding kernel and incidence identities.
Equations (AP.35)--(AP.36) are the exact map to the source's
orbit-factorial presentation.  The ungrouped matrices in
(AP.26)--(AP.34) continue to retain all their columns.

All zero-rank cases follow without inversion of \(\mathsf A_s\)
or \(\mathsf B_s\).  If their square-root product vanishes,
then \(p_M^s=0\), \(W_M^s=0\), and
\(\epsilon_M^s=0\), even if the kernel update is nonzero.
If \(\pi_M^s=1\), (AP.32) forces \(U_s=0\), with
the same conclusion.  The dimension \(N_s\) is always at
least one; if it is one, a Hermitian involution has only one
nonzero eigenspace, so one of \(U_s,V_s\) vanishes and
the weight is zero.  This argument also covers a zero-dimensional
reflection eigenspace in a larger target.  It neither removes the
retained quotient nor asserts existence of a particular actual packet.
For \(k=1\), every orbit cardinality is one and \(I=L=I_d\);
all the displayed formulas recover (AP.15)--(AP.21) exactly.

Finally the symmetric source is still \(R_MI\), and its two
actual boundary maps are \(Y_MI\) and \(C_MI\) in
(SS.22).  The original relation-graph map, followed by the bounded
permutation projection on its Hilbert target, realizes the redundant
frontier used here.  Grouping it by (AP.35) gives the invariant
relation layer of dimension \(p_{\le k}(M+1)\), as proved
in (SS.22)--(SS.24).  Its primitive is the original fixed-order
division primitive followed by \(\mathsf S_k\); the cochain
identity in (SS.3)--(SS.4) preserves its boundary and all signs.
At invariant support fibres the same maps lift with that support;
vanishing amplitudes remain the corresponding supported zero.
For other masks the original orbit-joined transport and its kernel
remain as stated after (SS.9).  Thus (AP.30) is a proved bound
for the same finite arithmetic cohomology summand, with no assumed
asymptotic estimate for either \(\Gamma_{k,M}\) or
\(\pi_M^s\).


Section~\ref{sec:quartet-dissipation-continuation} carries these
same two parities through the original graph frame into the actual
derivative layer. The equality
\(-\widehat Y^*\widehat C=UV^*\) proved there supplies the
exact original derivative cost for the positive operator here.
Sections~\ref{sec:theta-stieltjes-pair} and
\ref{sec:spectral-sum-stieltjes-pair} calculate its one-factor and
symmetric-square counterparts through explicit quadratic algebra
maps, with all local nilpotents and both reflection components.
